\section{An equivariant map with additive noise}\label{sec:equivariant-noise}

We prove Lemma~\ref{lem:equivariant-noise}, whose equivariant map supplies the
completeness labeling of the matrix game and whose linear detection bound
supplies its spectral comparison.  The construction separates a nonlinear
observable from families of linear observables under the same additive noise.  All vector spaces and
duals in this section are over $\F_2$, unless a different field is indicated.
Restriction of a subspace of a dual means its image under the restriction map.

We first construct a gadget with a fixed smaller output space, in which
shifts induce output translations.  After identifying shifts that induce
the same translation, this gadget will detect every linear family whose
restrictions fill the dual of the output space.  We then combine copies
indexed by injections of that space into $K$: a sufficiently high-rank
family on $K$ restricts onto the full output dual on most copies.  This second step allows every sufficiently large
dimension of $K$ while keeping the rank threshold fixed.

For the recursive maps, Lemma~\ref{lem:equivariant-noise-noise} computes the
nonlinear change probability, while
Proposition~\ref{prop:equivariant-noise-all-extensions-detection} bounds detection by every
linear family extending the pullbacks of all output functionals at the chosen height.  Section~\ref{subsec:equivariant-noise-enlargement}
passes both conclusions through the quotient and then enlarges the output
space; it contains the deduction of Lemma~\ref{lem:equivariant-noise}.

The fixed-output construction begins with a quadratic block.  On an
independent uniform input, a uniform nonzero perturbation in one of its
update spaces can leave the nonlinear output unchanged.  For most
subspaces of linear observations, restriction to the same update space
loses at most one dimension, and the loss probability is bounded by
three times the rank times this nonlinear stability probability.
The sum $\sum_{j=1}^r1/j$ in Equation~\eqref{eq:equivariant-noise-subspace-conditions-rank-loss}
removes the rank factor from the expected loss per level.  Random changes of coordinates
make it unlikely that a descending family contains distinct nonzero elements
proportional over the field or, for some binary linear map $g$, more than
three times its dimension nonzero solutions $z$ of
$z\cdot g(z)=\sqrt{z_1z_2z_3}$, even when later linear maps
depend on earlier choices.
Concatenation therefore makes the nonlinear change probability small while
leaving a constant probability of positive rank.  At the final leaf, a
surviving family detects an independent uniform symbol with constant
probability.

\subsection{The quadratic block}

Fix an integer $r_0\ge1$, whose size will eventually depend only on $p_*$.
Let
$F=\F_{2^d}$, $q=2^d$, and regard $X=B=F^3$ as binary spaces, each of
dimension $b=3d$.  The integer $d$ will be chosen sufficiently large in terms
of $r_0$.

Use the trace pairing to identify both $X^*$ and $B^*$ with $F^3$:
the vector $z\in F^3$ represents the linear functional
$y\mapsto\operatorname{Tr}_{F/\F_2}(z\cdot y)$, where the dot denotes the
field dot product.  We abbreviate the trace by $\operatorname{Tr}$.
The trace pairing is nondegenerate, and
$\operatorname{Tr}(t^2)=\operatorname{Tr}(t)$.  Every element of $F$ has a
unique square root.

Define
\begin{equation}\label{eq:equivariant-noise-quadratic}
 Q(x_1,x_2,x_3)=(x_2x_3,x_1x_3,x_1x_2),
 \qquad
 D(x,a)=Q(x+a)+Q(x)+Q(a).
\end{equation}
The block observable is
\[
 C_{\mathrm{blk}}:X\oplus B\longrightarrow B,
 \qquad C_{\mathrm{blk}}(x,y)=y+Q(x).
\]
It is equivariant under the vertical copy of $B$:
$C_{\mathrm{blk}}(x,y+b')=C_{\mathrm{blk}}(x,y)+b'$ for $b'\in B$.
For $z\in B^*$, the scalar observation $z\circ C_{\mathrm{blk}}$ changes
by $z(b')$ under this shift.  We shall compare it with linear functionals
whose restriction to the vertical $B$ is also $z$.

For a one-dimensional $F$-subspace $A=Fv$ of $F^3$, put
\begin{equation}\label{eq:equivariant-noise-block-subspaces}
 D_A=\{y\in B:v\cdot y=0\},
 \qquad
 U_A=\{(a,Q(a)+w):a\in A,\ w\in D_A\}\subseteq X\oplus B.
\end{equation}
These definitions are independent of the chosen nonzero generator $v$.
There are $q^2+q+1$ such lines; set
\begin{equation}\label{eq:equivariant-noise-theta}
 \theta=\frac{1}{q^2+q+1}.
\end{equation}

\begin{lemma}\label{lem:equivariant-noise-block}
For each one-dimensional $F$-subspace $A\le F^3$, the set $U_A$
in Equation~\eqref{eq:equivariant-noise-block-subspaces} is a binary subspace of
dimension $b$, and the trace annihilator of $D_A$ in $B^*=F^3$ is $A$.
Moreover,
\[
 D(x,a)\in D_A\qquad(x\in X,\ a\in A).
\]
As $A$ ranges over these lines, the subspaces $U_A$ span $X\oplus B$.
\end{lemma}

\begin{proof}
The three coordinates of $D(x,a)$ are
\[
 x_2a_3+a_2x_3,\qquad
 x_1a_3+a_1x_3,\qquad
 x_1a_2+a_1x_2.
\]
When $a=tv$, every term in $v\cdot D(x,a)$ occurs twice, so their sum is
zero.  When both arguments are multiples of $v$, each coordinate of the
polar expression is zero.  Thus $Q|_A$ is binary additive, and the
parametrization of $U_A$ by $A\oplus D_A$ is a binary linear bijection.
Its dimension is $d+2d=b$.

Every $z\in A$ annihilates $D_A$ under the trace pairing.  The annihilator
has binary dimension $3d-2d=d$, so it is exactly $A$.
Finally, for each coordinate line $A=Fe_i$, the space $U_A$ contains
$(Fe_i,0)$, because $Q(te_i)=0$, as well as $\{0\}\oplus D_A$.
These three coordinate lines and their vertical hyperplanes span
$X\oplus B$.
\end{proof}

\paragraph{The nonlinear response.}
The form of $U_A$ makes the quadratic term of an update cancel.  For
$(a,Q(a)+w)\in U_A$,
\[
 C_{\mathrm{blk}}\bigl((x,y)+(a,Q(a)+w)\bigr)
       +C_{\mathrm{blk}}(x,y)
 =w+D(x,a).
\]
The right-hand side lies in $D_A$ by Lemma~\ref{lem:equivariant-noise-block}.
The following calculation gives the exact law that the recursion will use.

\begin{lemma}[Response to one block update]\label{lem:equivariant-noise-block-noise}
For every one-dimensional $F$-subspace $A\le F^3$ and independent
uniform $(x,y)\in X\oplus B$ and $\zeta\in U_A\setminus\{0\}$,
\[
 \Pr\bigl[C_{\mathrm{blk}}((x,y)+\zeta)
             \ne C_{\mathrm{blk}}(x,y)\bigr]=1-\theta.
\]
\end{lemma}

\begin{proof}
First take $\zeta$ uniform in all of $U_A$, still independently of the
uniform input $(x,y)$.  Its parametrization
$(a,Q(a)+w)$ has independent uniform $a\in A$ and $w\in D_A$.
For every fixed $x,a$, the output difference $w+D(x,a)$ is uniform in
$D_A$, so it is nonzero with probability $1-q^{-2}$.  The zero
perturbation has mass $q^{-3}$ and never changes the output.  Conditioning
on a nonzero perturbation therefore gives
\[
 \frac{1-q^{-2}}{1-q^{-3}}
 =\frac{q(q+1)}{q^2+q+1}=1-\theta.
\]
\end{proof}

The lemma concerns an independent uniform input $(x,y)$ and a uniform
nonzero update.  In the recursion constructed below, $(x,y)$ will be the
aggregate of the embedded child outputs and the update will come from a
lower level; we will verify this conditional law before applying the
lemma.

\paragraph{Linear restrictions.}
Fix a binary subspace $S\le B^*$.  A linear functional on $X\oplus B$
whose restriction to the vertical $B$ is $z\in S$ has a unique form
$(\gamma_z,z)\in X^*\oplus B^*$.  Thus a binary-linear choice of one such
extension for each $z\in S$ is specified by a binary linear map
$\gamma:S\to X^*$.  We ask which of these chosen extensions vanish when
restricted to $U_A$.  The restriction loses rank exactly when its kernel
contains a nonzero element.

\begin{lemma}[Kernel of a block restriction]\label{lem:equivariant-noise-block-kernel}
Let $S\le B^*$ and let $\gamma:S\to X^*$ be binary linear.  For each line
$A$, define
\[
 \kappa_{A,\gamma}:S\longrightarrow U_A^*,
 \qquad \kappa_{A,\gamma}(z)=(\gamma_z,z)|_{U_A}.
\]
Let $g:S\to F^3$ represent $\gamma$ under the trace identification:
$\gamma_z(x)=\operatorname{Tr}(g(z)\cdot x)$ for $z\in S$, $x\in X$.
Then, for every nonzero $z\in S$,
\[
 \kappa_{A,\gamma}(z)=0
 \quad\Longleftrightarrow\quad
 A=Fz
 \quad\text{and}\quad
 z\cdot g(z)=\sqrt{z_1z_2z_3}.
\]
\end{lemma}

\begin{proof}
If $z$ lies in the kernel, it annihilates every vertical vector
$(0,w)$ with $w\in D_A$.  Lemma~\ref{lem:equivariant-noise-block} therefore gives
$z\in A$, and for $z\ne0$ this means $A=Fz$.  On the remaining vectors
$(tz,Q(tz))$, $t\in F$, the functional has value
\begin{equation}\label{eq:equivariant-noise-kernel-trace}
 \operatorname{Tr}\bigl(t\,g(z)\cdot z+t^2z_1z_2z_3\bigr).
\end{equation}
The three summands of $z\cdot Q(tz)$ are equal, and three equals one in
characteristic two.  Also
\[
 \operatorname{Tr}(t^2c)=\operatorname{Tr}(t\sqrt c).
\]
Nondegeneracy of the trace pairing shows that
Equation~\eqref{eq:equivariant-noise-kernel-trace} vanishes for every $t$ exactly when
$z\cdot g(z)=\sqrt{z_1z_2z_3}$.  Together with $z\in A$, this condition
is sufficient as well: the vertical part is already annihilated.
\end{proof}

\subsection{Two conditions on spaces of linear functionals}

The kernel criterion suggests two sufficient controls.  We shall use a
family in which no field line contains two nonzero linear functionals and only a few
linear functionals satisfy the equation $z\cdot g(z)=\sqrt{z_1z_2z_3}$ for any one choice of extensions.
The uniformity over that choice will later allow a family of extensions to depend on earlier
random changes of coordinates.

\begin{definition}
For $r\ge1$, let $\mathcal H_r$ consist of the binary subspaces
$S\le F^3$ of dimension $r$ such that no two distinct nonzero elements of $S$ are proportional over $F$, and,
for every binary linear map $g:S\to F^3$, at most $3r$ nonzero $z\in S$
satisfy
\begin{equation}\label{eq:equivariant-noise-kernel-equation}
 z\cdot g(z)=\sqrt{z_1z_2z_3}.
\end{equation}
\end{definition}

Write $H_0=0$ and $H_r=\sum_{j=1}^r1/j$ for positive integers $r$.
Fix one map $\gamma:S\to X^*$ before choosing a uniform line $A$, and let
$r_A=\dim\im\kappa_{A,\gamma}$, where $r=\dim S>0$.
Lemma~\ref{lem:equivariant-noise-block-kernel} now bounds the loss.  If $S\in\mathcal H_r$,
each field line contains at most one nonzero element of $S$, so the kernel
has dimension at most one.  For the fixed map $g$ associated with
$\gamma$, at most $3r$ linear functionals satisfy Equation~\eqref{eq:equivariant-noise-kernel-equation};
hence at most $3r$ lines can cause any rank loss.  Consequently
\begin{equation}\label{eq:equivariant-noise-subspace-conditions-rank-loss}
 r-r_A\le1,\qquad
 \Pr_A[r_A<r]\le3r\theta,\qquad
 \E_A[H_r-H_{r_A}]\le3\theta.
\end{equation}
The last bound uses $H_r-H_{r-1}=1/r$.
For an arbitrary rank-$r$ space, a loss still requires $A=Fz$ for some
nonzero $z\in S$.  Therefore
\begin{equation}\label{eq:equivariant-noise-arbitrary-rank-loss}
 \Pr_A[r_A<r]\le(2^r-1)\theta\le2^r\theta,
 \qquad H_r-H_{r_A}\le H_r.
\end{equation}
These estimates concern rank under restriction.  Detection of the eventual
noise will follow only after a positive-rank family reaches a leaf.

It remains to show that subspaces in $\mathcal H_r$ are abundant.

\begin{lemma}\label{lem:equivariant-noise-subspace-conditions}
Fix an integer $r\ge1$.  As $d\to\infty$, with $F=\F_{2^d}$,
the probability that a uniform binary $r$-dimensional subspace of $F^3$
belongs to $\mathcal H_r$ tends to one.
\end{lemma}

\begin{proof}
Choose independent uniform $x_{ij}\in F$, for $1\le i\le3$ and
$1\le j\le r$, and parametrize a binary linear map by
\[
 z_i(t)=\sum_{j=1}^r x_{ij}t_j,
 \qquad t\in\F_2^r.
\]
Each nonzero input has uniform image in $F^3$, so the probability that this
map is not injective is at most $(2^r-1)q^{-3}$.  Conditioned on
injectivity, its image is a uniform binary $r$-subspace: every such
subspace has the same number of ordered binary bases.  Distinct nonzero
binary inputs are linearly independent also over $F$, so their images are
independent uniform elements of $F^3$.  A union bound therefore shows that
the nonproportionality condition fails with probability $O_r(q^{-2})$,
apart from the event of noninjectivity.

To prove the bound on the number of solutions simultaneously for all $g$, we eliminate
the unknown basis values of $g$ by a column-span condition.  A nonzero
minor will then bound its probability without taking a union over the
field-dependent set of maps $g$.
Fix a set $E$ of more than $3r$ distinct nonzero inputs in $\F_2^r$.
On an injective parametrization, the values of a binary linear $g$ on a
basis can be arbitrary vectors in $F^3$.  Thus the equality on all of $E$
implies that the vector
\[
 c=\bigl(\sqrt{z_1(t)z_2(t)z_3(t)}\bigr)_{t\in E}
\]
lies in the $F$-column span of the matrix with $3r$ columns
\[
 M_{t,(i,j)}=t_jz_i(t),
 \qquad t\in E,\quad 1\le i\le3,\quad 1\le j\le r.
\]
Conversely, membership supplies the basis values of such a $g$.
In particular, this column-membership condition already accounts for all
possible maps $g$.

First regard the $x_{ij}$ as independent indeterminates and set
$L=\F_2(x_{ij})$.  We claim that $c$ is not in the column span of $M$ even
after extending scalars to an algebraic closure $\overline L$.
The following $r^3$ elements are linearly independent over $L$:
\[
 \sqrt{x_{1j}x_{2k}x_{3l}},
 \qquad 1\le j,k,l\le r.
\]
Indeed, after clearing denominators, a dependence would be a polynomial
identity in the square roots of the
indeterminates.  The displayed monomials have distinct parity vectors,
whereas coefficients polynomial in the $x_{ij}$ have even exponents in
those square roots.  Terms from different parity vectors cannot cancel.

The expansion of $c$ in these independent elements has coefficient vectors
\[
 v_{jkl}=(t_jt_kt_l)_{t\in E}.
\]
Suppose that $c$ belonged to the column span over $\overline L$.
Every row vector over $L$ annihilating $M$ would then annihilate $c$.
The independence just proved implies that it would annihilate every
$v_{jkl}$.  By double annihilation in the finite-dimensional space
$L^E$, every $v_{jkl}$ would belong to the column span over $L$ itself.

For $n\ge1$, let $V_n\le L^E$ be the evaluation span of all
positive-degree monomials in the binary coordinates $t_1,\ldots,t_r$
of degree at most $n$.  Repeated indices in $v_{jkl}$ give all monomials
of degrees one and two as well as three, because $t_i^2=t_i$ on the
evaluation set.  Each column $t\mapsto t_jz_i(t)=\sum_kx_{ik}t_jt_k$ of
$M$ is an $L$-combination of such monomials of degree at most two.  The
preceding conclusion and this form of the columns therefore imply
\[
 V_3\subseteq\operatorname{col}_L(M)\subseteq V_2\subseteq V_3.
\]
These three spaces are equal.  Multiplication by any $t_i$ sends $V_2$
into $V_3$, so preserves their common value.  Since this space contains
the degree-one monomials, it contains evaluations of every positive-degree
monomial.  These evaluations span all of $L^E$: for any nonzero binary
point, its usual indicator polynomial is a product of factors $t_i$ and
$1+t_i$ with at least one factor $t_i$, and thus has no constant monomial.
Its restriction to $E$ gives the corresponding point indicator.
Consequently the column span would have dimension $|E|>3r$, a
contradiction.

To pass from this algebraic statement to finite fields, square the entries
of both $M$ and $c$.  The resulting matrices have polynomial entries,
and the appended column has entries $z_1(t)z_2(t)z_3(t)$.
Let $h=\rank_L M\le3r$.  Squaring preserves rank over $\overline L$,
so the augmented squared matrix has a nonzero minor of size $h+1$.
This minor is a nonzero polynomial of degree at most $2h+3\le6r+3$.
Indeed, the inputs $t\in E$ are fixed: a squared entry of $M$ is
quadratic in the indeterminates $x_{ij}$, and an appended entry is
cubic.  Each determinant term uses one appended entry and $h$ other
entries.  The nonzero polynomial has coefficients in $\F_2$, so remains
nonzero over every extension $\F_{2^d}$.
All unaugmented minors of size $h+1$ vanish identically, so the
unaugmented rank cannot increase on specialization.  It follows that
specialized column membership forces this particular nonzero polynomial
to vanish.  For uniform independent $x_{ij}\in F$, this event has
probability at most $(6r+3)/q$ by the polynomial zero bound~\cite[Corollary~1]{Schwartz80};
see also the earlier sparse-polynomial formulation in~\cite[Section~3.1]{Zippel79}.

If the bound on the number of solutions fails, it fails on some set $E$ of size
$3r+1$.  There are finitely many such sets depending only on $r$; if
$2^r-1\le3r$ there are none and the condition is automatic.  A union
bound, followed by conditioning on injectivity, proves the lemma.
\end{proof}

The defining conditions of $\mathcal H_r$ therefore control every possible map $g$.  We now fix a
field where subspaces outside $\mathcal H_r$ are rare at every rank $r\le r_0$.  Combining the two bounds in
Equation~\eqref{eq:equivariant-noise-arbitrary-rank-loss} bounds the expected decrease in $H_r$
over a uniform line by $H_{r_0}2^{r_0}\theta$ at ranks at most $r_0$.
To limit the contribution from subspaces outside $\mathcal H_r$ of probability at most
$\eps_0$ to $\theta$, choose a positive $\eps_0$ with
\begin{equation}\label{eq:equivariant-noise-subspace-conditions-error}
 H_{r_0}2^{r_0}\eps_0\le1,
 \qquad \eps_0<1.
\end{equation}
By Lemma~\ref{lem:equivariant-noise-subspace-conditions}, we may fix $d$ sufficiently large that
$b=3d\ge r_0$ and, for every $1\le r\le r_0$, a uniform rank-$r$
subspace of $B^*$ lies outside $\mathcal H_r$ with probability at most $\eps_0$.
In particular, there exists a subspace in $\mathcal H_{r_0}$, which we fix
once and for all.

\subsection{Concatenation and its noise}

We keep the output space $B$ fixed while concatenating the quadratic blocks.
One noisy branch will be chosen at each level.  The next lemma computes how
fast the nonlinear change probability decays; we will then track the rank
of linear restrictions along the same branch.

We construct binary input spaces $P_t$, shift subspaces $R_t\le P_t$,
linear surjections $\lambda_t:R_t\to B$, and maps $C_t:P_t\to B$ such that
\begin{equation}\label{eq:equivariant-noise-recursive-equivariance}
 C_t(u+h)=C_t(u)+\lambda_t(h)
 \qquad(h\in R_t).
\end{equation}
A shift $h\in R_t$ thus represents the output translation $\lambda_t(h)$;
different shifts may represent the same translation.  This form of the
shift action is preserved by aggregation.  At the end, quotienting by the
shifts of value zero will turn the output translations into an actual
subspace of the input.

At height zero take $P_0=R_0=B$ and $C_0=\lambda_0=\operatorname{id}_B$.
Let $\mathcal I$ consist of all pairs $(A,J)$, where $A$ is a line in
$F^3$ and $J:B\to U_A$ is a binary linear isomorphism.  Each line has
the same number of such isomorphisms.

We include one child for every $(A,J)$ in a single fixed parent map.
Choosing a uniform noisy child will then sample a uniform line and a fresh
uniform isomorphism.
For $t\ge1$, take $P_t=\prod_{(A,J)\in\mathcal I}P_{t-1}$ and set
\begin{equation}\label{eq:equivariant-noise-recursive-output}
 (x,y)=\sum_{A,J}J C_{t-1}(u_{AJ}),
 \qquad C_t(u)=C_{\mathrm{blk}}(x,y)=y+Q(x).
\end{equation}
A tuple of child shifts changes this aggregate by the sum of their embedded
output translations.  To make the parent change a pure output translation,
we require the $X$-component of this sum to vanish.  Accordingly, define
\begin{equation}\label{eq:equivariant-noise-recursive-shifts}
 R_t=\left\{(h_{AJ})\in\prod_{A,J}R_{t-1}:
       \sum_{A,J}J\lambda_{t-1}(h_{AJ})\in\{0\}\oplus B\right\},
\end{equation}
and let $\lambda_t(h)=b'$ when this aggregate is $(0,b')$.
These are linear definitions, and Equation~\eqref{eq:equivariant-noise-recursive-equivariance}
follows immediately from Equation~\eqref{eq:equivariant-noise-recursive-output}.
Surjectivity follows inductively from Lemma~\ref{lem:equivariant-noise-block}: the
child output translations can be chosen arbitrarily, and their aggregate map
is onto $X\oplus B$.
Equation~\eqref{eq:equivariant-noise-recursive-equivariance} and surjectivity also show
that $C_t$ sends a uniform input to a uniform element of $B$.

Let $\xi_t$ be the following additive noise in $P_t$.
At height zero it is uniform in $B$.  At a positive height, choose a
uniform child $(A,J)$ and embed an independent copy of $\xi_{t-1}$ in
that child, with zero in all other children.  Equivalently, choose a
uniform leaf in the product tree and add a uniform symbol of $B$ at that
leaf.  Each path choice can be sampled as a uniform $A$ followed by a
uniform $J$, independently at successive levels.  All these distributions
are rational. Figure~\ref{fig:equivariant-noise-recursion} illustrates one level
of the recursion and the selected noisy branch.
\begin{figure}[!htbp]
\centering
\begin{tikzpicture}[
 every node/.style={font=\small,align=center},
 block/.style={draw=linkblue!80!black,rounded corners=2pt,
 text width=3.3cm,inner sep=7pt,minimum height=14mm},
 flow/.style={-{Latex[length=2mm]},thick}]
\node[block] (left) at (0,0)
 {One child input\\$u_{AJ}\in P_{t-1}$\\$J C_{t-1}(u_{AJ})$};
\node[block,draw=red!70!black] (selected) at (4.8,0)
 {Selected child input\\$u_{A_*J_*}+\xi_{t-1}$\\$J_*C_{t-1}(u_{A_*J_*}+\xi_{t-1})$};
\node[block] (right) at (9.6,0)
 {Other child inputs\\$u_{AJ}\in P_{t-1}$\\$J C_{t-1}(u_{AJ})$};
\node at (2.4,0) {$\cdots$};
\node at (7.2,0) {$\cdots$};
\node[block,text width=7cm] (sum) at (4.8,-2)
 {Sum over every child $(A,J)\in\mathcal I$\\
 noisy aggregate $(x',y')=(x,y)+J_*\Delta$};
\draw[flow] (left.south) -- (sum.north west);
\draw[flow,draw=red!70!black] (selected) -- (sum);
\draw[flow] (right.south) -- (sum.north east);
\node[block,text width=7cm] (out) at (4.8,-3.8)
 {Noisy parent output $y'+Q(x')\in B$};
\draw[flow] (sum) -- (out);
\end{tikzpicture}
\caption{One level of the recursion for the equivariant map, shown under its additive noise.
The input space is a product over all child indices $(A,J)$; only one
uniformly chosen child is perturbed. Each child output is embedded in
$U_A\le X\oplus B$, all embedded outputs are added, and the quadratic
map is applied once at the parent. Here $(x,y)$ is the initial aggregate,
and $\Delta$ is the selected child's output difference, so the noisy
aggregate is $(x',y')=(x,y)+J_*\Delta$. The omitted children obey the
same rule. The highlighted perturbation may disappear at the parent.}
\label{fig:equivariant-noise-recursion}
\end{figure}

\FloatBarrier

\begin{lemma}\label{lem:equivariant-noise-noise}
For every integer $t\ge0$, the recursive map $C_t$ and its noise
$\xi_t$ satisfy
\begin{equation}\label{eq:equivariant-noise-change-probability}
 \Pr[C_t(u+\xi_t)\ne C_t(u)]
     =(1-q^{-3})(1-\theta)^t,
\end{equation}
where $u$ is uniform in $P_t$ and independent of $\xi_t$.
\end{lemma}

\begin{proof}
Write $p_t=\Pr[C_t(u+\xi_t)\ne C_t(u)]$ under the independent
sampling in the statement.  The identity at height zero is immediate.
Consider one higher level and
condition on the selected child $(A,J)$.  Write $\Delta$ for its output
difference.  All child copies have the same function and the same inner
noise law, so the law of $\Delta$ is independent of the selected index.

We need also the conditional law of the child's initial output.
For any fixed child perturbation $a$, Equation~\eqref{eq:equivariant-noise-recursive-equivariance}
gives
\[
 C_{t-1}(u+a+h)+C_{t-1}(u+h)
 =C_{t-1}(u+a)+C_{t-1}(u)
 \qquad(h\in R_{t-1}).
\]
Thus the difference is invariant under these shifts, while the initial
output is translated by every element of $B$.  On uniform input its
initial output remains uniform conditional on the difference, even if the
inner perturbation is also specified.  All other initial child outputs
remain independently uniform.  Because the aggregate map is onto
$X\oplus B$, the initial aggregate $(x,y)$ is therefore uniform
conditional on $(A,J,\Delta)$.  In particular it is independent of the
aggregate perturbation $J\Delta$.

Conditional on $\Delta\ne0$ and $A$, the isomorphism $J$ remains uniform,
because the law of $\Delta$ is independent of the selected child index.
Thus $J\Delta$ is uniform in $U_A\setminus\{0\}$.  The conditional
uniformity just proved makes the initial aggregate independent of this
update.  Lemma~\ref{lem:equivariant-noise-block-noise} therefore gives parent change
probability $1-\theta$ conditional on a nonzero child difference.  When
$\Delta=0$, the aggregate does not change.  Consequently
$p_t=p_{t-1}(1-\theta)$, proving the identity.
\end{proof}

\subsection{Linear extensions and loss of rank}

The nonlinear change probability now decays with the height.  For the
fixed-$B$ gadget, our remaining target is detection whenever a family's
restrictions fill $B^*$.  We will prove this by tracking every binary-linear
choice of extensions of these output linear functionals through the shift action.

For $S\le B^*$, consider a binary linear family of extensions on $P_t$, that is, a map
$z\mapsto\psi_z$ from $S$ to $P_t^*$ satisfying
\begin{equation}\label{eq:equivariant-noise-extensions-definition}
 \psi_z|_{R_t}=z\circ\lambda_t.
\end{equation}
We next describe every such family at a parent.

\begin{lemma}\label{lem:equivariant-noise-extensions-factorization}
Let $t\ge1$ and let $z\mapsto\psi_z$ be a linear family of extensions on $P_t$ with
domain $S\le B^*$.  For each child $(A,J)$, restriction to its copy
of $R_{t-1}$ factors through $\lambda_{t-1}:R_{t-1}\to B$.
The resulting linear functionals on $B$ are
\begin{equation}\label{eq:equivariant-noise-child-functionals}
 \phi_{AJ}(z)=J^*\bigl((\gamma_z,z)|_{U_A}\bigr)\in B^*,
\end{equation}
where $\gamma:S\to X^*$ is a single binary linear map common to all the
children.
\end{lemma}

\begin{proof}
A shift supported on one child and lying in $\ker\lambda_{t-1}$ is a
parent shift with output translation zero.  Each $\psi_z$ vanishes on it,
proving the first assertion.

Let $T=\prod_{A,J}B$ be the space of child output translations, with aggregate
map $\mathcal A:T\to X\oplus B$.  This map is onto.  The sum of the
child linear functionals, minus $z$ applied to the second component of
$\mathcal A$, vanishes whenever the first component of $\mathcal A$ is
zero: every such tuple of translations can be supplied by child shifts,
giving a parent shift.  The difference therefore factors uniquely through the
first component of $\mathcal A$, which is onto $X$.  Call this linear functional
$\gamma_z$.  Uniqueness shows that $\gamma_z$ depends linearly on $z$.
Evaluating on a tuple supported at one child gives
Equation~\eqref{eq:equivariant-noise-child-functionals}.
\end{proof}

After choosing a child, let $S'=\im\phi_{AJ}$ and choose a linear section
$\iota:S'\to S$ of $\phi_{AJ}$.  Restricting $\psi_{\iota(z')}$ to that
child input gives a linear family of extensions with domain $S'$ on $P_{t-1}$.
In fact, if $\jmath_{AJ}$ embeds that child in the parent product, then
for $h\in R_{t-1}$,
\[
 (\psi_{\iota(z')}\circ\jmath_{AJ})(h)
 =\phi_{AJ}(\iota(z'))(\lambda_{t-1}h)
 =z'(\lambda_{t-1}h).
\]
This checks the shift action required of the new family of extensions even when
$\phi_{AJ}$ has a kernel.  It also identifies the surviving linear functional
as the restriction of an original linear functional, which is needed when
detection is passed back up the selected branch.
We may fix deterministic rules for all such sections.  In particular they
depend only on the already chosen path, and never on a future path index
or on the final noise symbol.  On descending the tree, every surviving
functional is a restriction of a member of the original linear family.

The factorization now connects the recursion to the local restriction
problem.  For a fixed parent family of extensions, its single map $\gamma$ is common to
every outgoing line $A$, and
$\phi_{AJ}=J^*\kappa_{A,\gamma}$.  Since $J^*$ is an isomorphism, it does
not change the kernel or rank.  We may therefore apply
Equations~\eqref{eq:equivariant-noise-subspace-conditions-rank-loss}
and~\eqref{eq:equivariant-noise-arbitrary-rank-loss} when the outgoing line is chosen
uniformly.  The next proof shows that the incoming isomorphisms make
membership in $\mathcal H_r$ typical at each positive subsequent rank $r$, even though the family of extensions
itself depends on the earlier path.

\begin{proposition}\label{prop:equivariant-noise-all-extensions-detection}
For the field and the bound $\eps_0$ on failure of the two conditions, fixed in
Equation~\eqref{eq:equivariant-noise-subspace-conditions-error} and the paragraph following it, put
\begin{equation}\label{eq:equivariant-noise-depth}
 s=\left\lfloor\frac{H_{r_0}}{8\theta}\right\rfloor.
\end{equation}
Every linear family of extensions $z\mapsto\psi_z$ on $P_s$ with domain $B^*$ satisfies
\begin{equation}\label{eq:equivariant-noise-all-extensions-detection}
 \Pr[\text{there exists }z\in B^*\text{ with }\psi_z(\xi_s)\ne0]
 \ge\frac14.
\end{equation}
\end{proposition}

\begin{proof}
Restrict the family with domain $B^*$ at the top to the fixed subspace in $\mathcal H_{r_0}$.  Follow the random leaf path, replacing the family at each
chosen child by the family of extensions on the image just described.  Let $S_j$ be its domain
after $j$ steps, and let $r_j=\dim S_j$, so initially $r_j=r_0$ at $j=0$.
Continue with the zero space after extinction.

We spell out the conditioning that permits the defining conditions of $\mathcal H_r$ to be used at every
level.  Let $\mathcal F_j$ be the information in all path choices before
the outgoing choice $(A_j,J_j)$ at step $j$.  The current space, family of extensions,
and map $\gamma_j$ are $\mathcal F_j$-measurable.
Conditional on $\mathcal F_j$ and $A_j$, the subspace
\[
 E_{j,A_j}
 =\im\bigl[z\mapsto(\gamma_{j,z},z)|_{U_{A_j}}\bigr]
 \le U_{A_j}^*
\]
is fixed.  The next domain and its rank are
\[
 S_{j+1}=J_j^*E_{j,A_j},
 \qquad r_{j+1}=\dim E_{j,A_j}.
\]
The rank is fixed before $J_j$ is revealed.  A uniform binary
isomorphism $J_j:B\to U_{A_j}$ makes $S_{j+1}$ a uniformly random
subspace of $B^*$ of that rank.  Therefore, if $B_j$ is the indicator
that $r_j>0$ and $S_j\notin\mathcal H_{r_j}$, then
\begin{equation}\label{eq:equivariant-noise-fresh-orientation}
 \E[B_{j+1}\mid\mathcal F_j,A_j]\le\eps_0.
\end{equation}
The next family of extensions may depend on $J_j$.  This causes no restriction on the
argument: membership in $\mathcal H_{r_{j+1}}$ bounds every binary linear map $g$ on the selected
subspace, including the map supplied by that family.

Let $\Delta_j=H_{r_j}-H_{r_{j+1}}\ge0$.
Conditional on $\mathcal F_j$, the outgoing $A_j$ is independent and
uniform, and the rank-loss estimates above apply to the fixed current
family of extensions.  Consequently
\begin{equation}\label{eq:equivariant-noise-conditional-sum-drop}
 \E[\Delta_j\mid\mathcal F_j]
 \le3\theta+H_{r_0}2^{r_0}\theta B_j.
\end{equation}
For $j=0$ the space belongs to $\mathcal H_{r_0}$, so its expected drop is at most
$3\theta$.  For $j\ge1$, average Equation~\eqref{eq:equivariant-noise-conditional-sum-drop}
over the preceding isomorphism using
Equation~\eqref{eq:equivariant-noise-fresh-orientation}.  This gives the unconditional
bound
\[
 \E\Delta_j
 \le\bigl(3+H_{r_0}2^{r_0}\eps_0\bigr)\theta
 \le4\theta.
\]

Summing the expected drops and using Equation~\eqref{eq:equivariant-noise-depth},
\[
 \E[H_{r_0}-H_{r_s}]\le4s\theta\le\frac{H_{r_0}}2.
\]
Extinction requires a total drop of $H_{r_0}$, so
$\Pr[r_s=0]\le1/2$.  At a leaf, $P_0=R_0=B$ and its family of extensions consists of
the linear functionals themselves.  Conditional on any surviving rank $r_s>0$,
an independent uniform symbol in $B$ is detected with probability
$1-2^{-r_s}\ge1/2$.  Every surviving functional is a restriction of the
original family along the selected leaf, so such a detection implies a
detection by that family of the embedded noise $\xi_s$.  Multiplying the
two bounds proves Equation~\eqref{eq:equivariant-noise-all-extensions-detection}.
\end{proof}

\subsection{Quotients and enlargement}\label{subsec:equivariant-noise-enlargement}

The fixed-$B$ construction is complete once its shift action is realized
by an actual subspace.  Put
\[
 V_0=P_s/\ker\lambda_s,
\]
where the kernel is taken inside $R_s\le P_s$, and let
$q_0:P_s\to V_0$ be the quotient map.
Equation~\eqref{eq:equivariant-noise-recursive-equivariance} makes $C_s$ invariant under
this kernel, so it descends to a map $\overline C:V_0\to B$.
The quotient $R_s/\ker\lambda_s$ is identified with $B$ by $\lambda_s$,
giving an inclusion $i_B:B\to V_0$ such that
\[
 q_0(h)=i_B(\lambda_s h)\qquad(h\in R_s).
\]
We henceforth identify this subspace with $B$.  Its shift action satisfies
\[
 \overline C(x+b')=\overline C(x)+b'
 \qquad(b'\in B).
\]
Let $\nu$ be the pushforward of the noise law of $\xi_s$.
Uniform input on $P_s$ pushes to uniform input on $V_0$, so the output
change probability remains $(1-q^{-3})(1-\theta)^s$.

Every linear subspace $G_0\le V_0^*$ whose restriction is all of $B^*$
satisfies
\begin{equation}\label{eq:equivariant-noise-base-detection}
 \Pr_{a\sim\nu}[\text{there exists }g\in G_0\text{ with }g(a)\ne0]
 \ge\frac14.
\end{equation}
Indeed, choose a linear section $\sigma:B^*\to G_0$ of restriction and
define $\psi_z=q_0^*\sigma(z)=\sigma(z)\circ q_0$.  For $h\in R_s$,
\[
 \psi_z(h)=\sigma(z)(q_0h)=z(\lambda_s h).
\]
Thus this pullback is a linear family of extensions with domain $B^*$.  Moreover,
$\psi_z(\xi_s)=\sigma(z)(q_0\xi_s)$, so detection by the family of extensions is detection
by the chosen section on the pushed-forward noise.
Proposition~\ref{prop:equivariant-noise-all-extensions-detection} gives the bound, and
detection by the section implies detection by $G_0$.

\begin{proof}[Proof of Lemma~\ref{lem:equivariant-noise}]
Choose $r_0$ sufficiently large that
\[
 \exp(1-H_{r_0}/8)\le p_*.
\]
Then choose $\eps_0$ as in Equation~\eqref{eq:equivariant-noise-subspace-conditions-error} and
fix $d$, and hence $b=3d$ and $\theta$, using
Lemma~\ref{lem:equivariant-noise-subspace-conditions}.  Construct the height $s$ in
Equation~\eqref{eq:equivariant-noise-depth} and the quotient $(V_0,B,\overline C,\nu)$
above.  Lemma~\ref{lem:equivariant-noise-noise} gives
\begin{equation}\label{eq:equivariant-noise-small-base-change}
 (1-q^{-3})(1-\theta)^s\le(1-\theta)^s
 \le\exp(-\theta s)
 \le\exp(\theta-H_{r_0}/8)
 \le p_*.
\end{equation}
Set $r_*=b+2$.  All choices so far depend only on $p_*$.

We now enlarge the output space.  The point of using random injections
$B\to K$ is that a family of restriction rank at least $b+2$ on $K$
restricts onto all of $B^*$ along most injections.  A copy of the base
gadget for each injection will therefore turn its full-restriction
detection bound into the required high-rank bound.

Fix $\ell\ge b+2$ and write $K=\F_2^\ell$.
Let $\mathcal L$ be the finite set of all binary linear injections
$L:B\to K$.  Take a product $P=\prod_{L\in\mathcal L}V_0$ and define
\[
 \widetilde C((x_L)_L)=\sum_{L\in\mathcal L}L\overline C(x_L).
\]
The product of shift subspaces $T=\prod_L B\le P$ acts on this output
through the linear map
\[
 \Lambda:T\longrightarrow K,
 \qquad \Lambda((b_L)_L)=\sum_L Lb_L.
\]
This map is onto, because every nonzero vector of $K$ is in the image of
some injection $B\to K$.  Quotient by its kernel inside $T$:
\[
 \mathcal V=P/\ker\Lambda.
\]
The map $\widetilde C$ descends to $C:\mathcal V\to K$, and
$T/\ker\Lambda$ identifies with the required shift subspace $K$.
This proves Equation~\eqref{eq:equivariant-noise-equivariance}.

On $P$, choose a uniform $L\in\mathcal L$ and update only that copy by
an independent sample from $\nu$.  Let $\mu$ be the resulting law after
the quotient.  On uniform product input the output sum changes exactly
when the selected copy's output changes, because $L$ is injective.
Uniform input pushes to uniform input on the quotient, so the change
probability is still $(1-q^{-3})(1-\theta)^s$.  Equation~\eqref{eq:equivariant-noise-small-base-change}
proves Equation~\eqref{eq:equivariant-noise-small-change}.  Every noise distribution
used here is obtained from uniform finite choices by pushforward, and
therefore is rational.

It remains to establish detection for an arbitrary
$G\le\mathcal V^*$.  Let $S=G|_K\le K^*$, of dimension $r\ge r_*$.
Pull $G$ back to $P$ and then restrict to its copy of $V_0$ indexed by
$L$; denote the resulting subspace of $V_0^*$ by $G_L$.
Its restriction to that copy's shift subspace $B$ is exactly
\[
 G_L|_B=L^*S.
\]
To see this, a shift $b'\in B$ supported on copy $L$ represents the
shift $Lb'\in K$ in the quotient, so a linear functional restricting to
$z\in S$ evaluates there as $z(Lb')$.

The map $L^*:S\to B^*$ fails to be onto if and only if there is a
nonzero $b'\in B$ with $Lb'\in S^\perp$.  For each such fixed $b'$, a
uniform injection $L$ makes $Lb'$ uniform among the nonzero vectors of
$K$.  Hence
\[
 \Pr_L[Lb'\in S^\perp]
 =\frac{2^{\ell-r}-1}{2^\ell-1}
 \le2^{-r}.
\]
A union bound gives
\[
 \Pr_L[L^*S\ne B^*]
 \le(2^b-1)2^{-r}
 \le2^{b-r}\le\frac14.
\]
On each of the remaining copies,
Equation~\eqref{eq:equivariant-noise-base-detection} gives detection probability at
least $1/4$ by $G_L$, and thus by $G$ on the embedded noise.
Averaging over the selected copy yields detection probability at least
$\frac34\cdot\frac14=\frac{3}{16}\ge\frac18$, proving
Equation~\eqref{eq:equivariant-noise-detection}.
The construction works for every $\ell\ge b+2$, with the already fixed
$r_*=b+2$, as required.

For completeness, all choices are effective when $p_*$ is rational.
Choose an integer $m\ge0$ with $2^{-m}\le p_*$ and then choose $r_0$ with
$H_{r_0}\ge8(m+1)$; since $e^{-m}\le2^{-m}$, this ensures the initial
bound on $r_0$.  Choose rational $\eps_0$ satisfying
Equation~\eqref{eq:equivariant-noise-subspace-conditions-error}.  For successive integers $d$ with
$3d\ge r_0$, enumerate the binary subspaces of $\F_{2^d}^3$ and the linear
maps in the definition of $\mathcal H_r$.  This tests membership exactly and
computes the fraction of subspaces outside $\mathcal H_r$ at each rank $r=1,\ldots,r_0$.
Lemma~\ref{lem:equivariant-noise-subspace-conditions} ensures that eventually all these fractions
are at most $\eps_0$.  Fix the first such field and a subspace in $\mathcal H_{r_0}$.  The remaining products, isomorphisms, injections, quotients, function
tables, and rational noise laws are obtained by finite enumeration and
binary linear algebra.  This gives the asserted terminating construction.
\end{proof}
